\Article{17}
\begin{center}\emph{Theorie des biquadratischen Restes} $1+i$.\end{center}
Der Modulus soll mit dem Reste keinen Theiler gemein haben, wir nehmen also an, dass von den Zahlen $a$ und $b$ die eine gerade, die andere ungerade sei. Die Vertheilung der eigentlichen Primitivreste in die vier Classen stellt folgendes Schema vor
\begin{center}\resizebox{.90\textwidth}{!}{\input{diagrams/native328_classes.tex}}\end{center}

Zu $n$ sind zu rechnen alle Zahlen auf
\[
\text{der Linie }0\ldots\tfrac12m\qquad\text{Anzahl }=g
\]
Zu $n'$ alle Zahlen auf
\[
\text{der Linie }0\ldots(\tfrac14+\tfrac14i)m\qquad\text{Anzahl }=g'
\]
Zu $n''$ alle Zahlen innerhalb
\[
\text{des Dreiecks }\tfrac12m,m,(\tfrac12+\tfrac12i)m\qquad\text{Anzahl }=h
\]
Zu $n'''$ alle Zahlen innerhalb
\[
\text{des Dreiecks }0,\tfrac12m,(\tfrac12+\tfrac12i)m\qquad\text{Anzahl }=h'
\]
und ausserdem alle Zahlen auf
\[
\text{der Linie }(\tfrac14+\tfrac14i)m\dotstwo(\tfrac12+\tfrac12i)m\qquad\text{Anzahl }=g''
\]

Man hat immer $g'+g''=g$, $aa+bb=p$, $\frac14(p-1)=g+g'+g''+h+h'$ Der Decident ist also
\SourceLeaf{329}

\[
D=S(0_{(+)},\tfrac12m,(\tfrac12+\tfrac12i)m_{(-)})+\tfrac12(p-1)+2g''
\]

Man nehme nun an, dass für den Modulus $m+1+i$
\[
\begin{array}{clc}
g&\text{übergehe in}&G\\
g'&&G'\\
g''&&G''\\
h&&H\\
h'&&H'
\end{array}
\]
so hat man
\[
\begin{aligned}
&\Delta S(0_{(+)},\tfrac12m,(\tfrac12+\tfrac12i)m_{(-)})\\
&\quad=+S(0_{(+)},(\tfrac12+\tfrac12i)m,(\tfrac12+\tfrac12i)m+i_{(-)})\\
&\qquad-S(0_{(+)},\tfrac12m,\tfrac12m+\tfrac12+\tfrac12i_{(-)})\\
&\qquad-S(\tfrac12m,(\tfrac12+\tfrac12i)m,(\tfrac12+\tfrac12i)m+i,\ \tfrac12m+\tfrac12+\tfrac12i)
\end{aligned}
\]

Das letzte dieser $S$ ist
\[
=\left[\tfrac12(a-b)\right]-\left[\tfrac12a\right]-S((\tfrac12+\tfrac12i)m,\tfrac12m,\tfrac12m+\tfrac12-\tfrac12i)
\]
wenn $a$ ungerade oder gerade
\[
\begin{aligned}
={}&\left[\tfrac12(a-b)\right]-\left[\tfrac12a\right]-S(-\tfrac12+\tfrac12i,\tfrac12m-\tfrac12+\tfrac12i,\tfrac12m+i)\\
={}&\left[\tfrac12(a-b)\right]-\left[\tfrac12a\right]+S(0_{(+)},\tfrac12m,\tfrac12m+\tfrac12+\tfrac12i_{(-)})\\
&-S(0_{(+)},(\tfrac12+\tfrac12i)m,(\tfrac12+\tfrac12i)m+i_{(-)})
\end{aligned}
\]

Also
\[
\begin{aligned}
\Delta D={}&-\left[\tfrac12(a-b)\right]+\left[\tfrac12a\right]\\
&+2S(0_{(+)},(\tfrac12+\tfrac12i)m,(\tfrac12+\tfrac12i)m+i_{(-)})+a+b+1\\
&-2S(0_{(+)},\tfrac12m,\tfrac12m+\tfrac12+\tfrac12i_{(-)})-2g''+2G''
\end{aligned}
\]

Die Bindezeichen gelten alle für den Fall, wo $a-b$ positiv ist, sonst nimmt man die entgegengesetzten.

Wir zerlegen ferner $S(0_{(+)},(\frac12+\frac12i)m,(\frac12+\frac12i)m+i_{(-)})$ in
\[
\begin{aligned}
&S(0_{(+)},(\tfrac14+\tfrac14i)m,(\tfrac14+\tfrac14m)+\tfrac12i_{(-)})\\
&+\left[\tfrac12(a-b)\right]-\left[\tfrac14(a-b)\right]\\
&-S((\tfrac12+\tfrac12i)m_{(-)},(\tfrac14+\tfrac14i)m,(\tfrac14+\tfrac14i)m-\tfrac12i_{(-)})
\end{aligned}
\]

\SourceLeaf{330}

Der letzte Theil
\[
\begin{aligned}
={}&-S(\tfrac12+\tfrac12i_{(-)},(-\tfrac14-\tfrac14i)m+\tfrac12+\tfrac12i,\ (-\tfrac14-\tfrac14i)m+\tfrac12_{(+)})\\
={}&-S(-\tfrac12-\tfrac12i_{(-)},(\tfrac14+\tfrac14i)m-\tfrac12-\tfrac12i,\ (\tfrac14+\tfrac14i)m-\tfrac12_{(+)})\\
={}&-S(-\tfrac12-\tfrac12i_{(+)},(\tfrac14+\tfrac14i)m-\tfrac12-\tfrac12i,\ (\tfrac14+\tfrac14i)m-\tfrac12_{(-)})+g''-G''\\
={}&S(0_{(+)},(\tfrac14+\tfrac14i)m,(\tfrac14+\tfrac14i)m+\tfrac12i_{(-)})\\
&-S(0_{(+)},\tfrac12m,\tfrac12m+\tfrac12+\tfrac12i_{(-)})+g''-G''
\end{aligned}
\]

Dadurch wird also
\[
\begin{aligned}
\Delta D={}&+\left[\tfrac12(a-b)\right]+\left[\tfrac12a\right]\\
&-4S(-\tfrac12-\tfrac12i_{(+)},\tfrac14+\tfrac14im-\tfrac12-\tfrac12i,\ (\tfrac14+\tfrac14i)m-\tfrac12_{(-)})\\
&-2\left[\tfrac14(a-b)\right]+a+b+1
\end{aligned}
\]

Für den Fall der Vermehrung des Modulus um $1-i,-1+i,-1-i$ ist keine besondere Untersuchung nöthig, weil offenbar die Moduli $m,im,-m,-im$ gleiche Decidenten haben. Wir haben also folgende Lehrsätze:

Ist der Decident des Modulus $a+bi,=D$, so sind die Decidenten von
{\small
\[
\begin{array}{l|l}
a+1+(b+1)i&D+a+b+1+[\tfrac12a]+[\tfrac12(a-b)]-2[\tfrac14(a-b)]\\
a+1+(b-1)i&D+a-b+1+[-\tfrac12b]+[-\tfrac12(a+b)]-2[-\tfrac14(a+b)]\\
a-1+(b+1)i&D-a+b+1+[\tfrac12b]+[\tfrac12(a+b)]-2[\tfrac14(a+b)]\\
a-1+(b-1)i&D-a-b+1+[-\tfrac12a]+[\tfrac12(b-a)]-2[\tfrac14(b-a)]
\end{array}
\]
}

Hieraus ferner
{\footnotesize
\[
\begin{array}{l|l|l}
&&\begin{gathered}\text{oder insofern }a\text{ un-}\\\text{gerade ist}\end{gathered}\\
a+2+bi&\begin{aligned}D-\tfrac{a-b-3}2+2[\tfrac12a]&-2[\tfrac{a-b}4]\\&-2[\tfrac{-a-b-2}4]\end{aligned}&D+\tfrac{b-a-1}2\\[1ex]
a-2+bi&\begin{aligned}D+\tfrac{a-b+3}2+2[-\tfrac12a]&-2[\tfrac{b-a}4]\\&-2[\tfrac{a+b-2}4]\end{aligned}&D\\[1ex]
a+(b+2)i&\begin{aligned}D-\tfrac{a+b-3}2+2[\tfrac12b]&-2[\tfrac{a+b}4]\\&-2[\tfrac{a-b-2}4]\end{aligned}&D+\tfrac{a+b-3}2\\[1ex]
a+(b-2)i&\begin{aligned}D+\tfrac{a+b+3}2+2[-\tfrac12b]&-2[\tfrac{-a-b}4]\\&-2[\tfrac{b-a-2}4]\end{aligned}&
\end{array}
\]
}

\[
\begin{array}{l|rl}
a+2+(b+2)i&2a-b+1+D&\text{oder }D+b-1\\
a+2+(b-2)i&-a-2b+1+D&D+a-1\\
a-2+(b+2)i&a+2b+1+D&D-a-1\\
a-2+(b-2)i&-2a+b+1+D&D-b-1
\end{array}
\]

\SourceLeaf{331}

\[
\begin{array}{l|l@{\qquad}l|l}
a+4+bi&D+a+b&a+8+bi&D+2\\
a+(b+4)i&D-a+b&&\\
a-4+bi&D-a-b&&\\
a+(b-4)i&D+a-b&&
\end{array}
\]
\[
\begin{array}{l|l}
a+4+(b+4)i&D+2b\\
a+4+(b-4)i&D+2a\\
a-4+(b+4)i&D+2a\\
a-4+(b-4)i&D+2b
\end{array}
\]

Das Resultat der vorhergehenden Untersuchungen ist also folgendes:
Für den Modulus $m=a+bi$, wo $a$ ungerade, $b$ gerade, wird
\[
\begin{aligned}
D\frac{1+i}{m}&=\tfrac18(-aa+2ab+bb-8b+1)\\
&\quad\text{(und wenn }a+bi=1+(2+2i)(\alpha+\beta i))\\
&\text{oder }\tfrac18(-aa+2ab-3bb+1)\equiv-(\alpha-\beta)^2-\beta
\end{aligned}
\]
\[
\begin{aligned}
D\frac{1-i}{m}&=\tfrac18(+aa+2ab-bb-8b-1)\\
&\text{oder }\tfrac18(+aa+2ab+3bb-1)
\end{aligned}
\]
\[
\begin{aligned}
D\frac{-1-i}{m}&=\tfrac18(-aa+2ab+bb+1)\\
&=\beta+\alpha\alpha+2\alpha\beta-\beta\beta\equiv-\beta+(\alpha+\beta)^2
\end{aligned}
\]
\[
\begin{aligned}
D\frac{-1+i}{m}&=\tfrac18(+aa+2ab-bb-1)\\
&=\alpha+\alpha\alpha-2\alpha\beta-\beta\beta\equiv-\alpha-(\alpha+\beta)^2
\end{aligned}
\]
\[
\begin{aligned}
D\frac2m&=\tfrac12ab\\
D\frac im&=\tfrac14(aa+bb-1)\\
D\frac{-1}m&=\tfrac12(aa+bb-1)
\end{aligned}
\]

\begin{center}\emph{Allgemeines Theorem über die Decidenten.}\end{center}
Es seien $A,B,C$ etc. ungleiche (unger. imag.) Primzahlen, deren keine die Zahl $M$ misst: alsdann ist
\[
D\frac{M}{A^\alpha B^\beta C^\gamma D^\delta}=\alpha D\frac MA+\beta D\frac MB+\gamma D\frac MC+\text{etc.}
\]
\[
M^{\frac14(aa+bb-1)}\equiv i^{D\frac{M}{a+bi}}\ (\mathrm{mod.}\ (a+bi)\quad\text{wenn }a+bi\text{ eine Primzahl}
\]

\[
\begin{aligned}
D\frac{1+i}m&=\tfrac18(-aa+2ab-3bb+1)\\
&=-\tfrac18(3(a-b)\mp1)(a-b\mp1)\quad\text{wenn }a\equiv\substack{1\\3}\\
D\frac{1-i}m&=\tfrac18(+aa+2ab+3bb-1)\\
D\frac{-1-i}m&=\tfrac18(-aa+2ab+bb+1)\\
&=\tfrac18(a-b\mp1)(a-b\mp3)\quad\text{wenn }a\equiv\pm1\\
D\frac{-1+i}m&=\tfrac18(+aa+2ab-bb-1)
\end{aligned}
\]

\SourceLeaf{332}

Allgemein $m\equiv1\ \mathrm{mod.}(1+i)$
\[
\begin{aligned}
D\frac{1+i}m&=-\text{P. Real. }\frac{(1+i)(m^4-1)}{16}=\text{Coëff. im. }\frac{m^4-1}{8+8i}\\
D\frac{1-i}m&=+\text{P. Real. }\frac{(1-i)(m^4-1)}{16}=\text{Coëff. im. }\frac{m^4-1}{8-8i}
\end{aligned}
\]

\[
\begin{array}{c|rrrrrrrr}
\multicolumn{9}{c}{1+i\ (\mathrm{mod.}\ 16)}\\\hline
a&\multicolumn{8}{l}{\quad b}\\
&0&2&4&6&8&10&12&14\\\hline
1&0&3&3&0&2&1&1&2\\
3&3&3&0&2&1&1&2&0\\
5&1&2&0&3&3&0&2&1\\
7&2&0&3&3&0&2&1&1\\
9&2&1&1&2&0&3&3&0\\
11&1&1&2&0&3&3&0&2\\
13&3&0&2&1&1&2&0&3\\
15&0&2&1&1&2&0&3&3
\end{array}
\]
